\documentclass[11pt,a4paper]{article}

\usepackage[utf8]{inputenc}
\usepackage[T1]{fontenc}
\usepackage[catalan]{babel}
\usepackage{lmodern}
\usepackage[a4paper,margin=2.5cm]{geometry}
\usepackage{enumitem}
\usepackage{hyperref}
\usepackage{parskip}

\setlist[itemize]{leftmargin=1.5em,itemsep=0.5em,topsep=0.5em}

% \title{\textbf{PRÀCTIQUES D'ESTADÍSTICA}\\[0.5em]
% \large Grau en Enginyeria Informàtica\\
% Curs 2026--27}
% \author{}
% \date{}

\begin{document}

\begin{center}
 {\textbf{PRÀCTIQUES D'ESTADÍSTICA}\\[0.5em]
\large Grau en Enginyeria Informàtica\\
Curs 2026--27}
\end{center}


\begin{itemize}
    \item Les pràctiques consten de 4 sessions ordinàries i dues proves d'avaluació a la 3a i 6a sessió programada, previstes segons el calendari (torn A) i els horaris ja establerts.

    \item Cada pràctica tindrà una durada de 1h i 40 minuts.

    \item L'assistència és obligatòria. Dues faltes injustificades suposarà la suspensió d'aquesta part de l'assignatura.

    \item L'objectiu de les sessions és, a través del treball autònom de l'estudiant, familiaritzar-se amb l'entorn i aprendre a resoldre problemes de probabilitat i estadística amb R. El professor farà els controls pertinents a l'aula i ajudarà a l'alumnat durant les sessions.

    A les dues proves d'avaluació, l'estudiant haurà de demostrar haver adquirit les competències bàsiques treballades durant les pràctiques.

    \item Si és possible, és altament recomanable que cada alumne dugui el seu propi ordinador. Evidentment, els ordinadors de l'Escola estan perquè els utilitzi qui ho vulgui i/o ho necessiti.
\end{itemize}

\subsubsection*{Grups i professorat}

\begin{tabular}{ll}
411: & Marc Magaña\\
412: & Juan J.~Rodríguez\\
413: & Juan J.~Rodríguez\\
414: & Juan J.~Rodríguez\\
415: &  Juan J.~Rodríguez\\
416: & Alan Morte\\
417: & Alan Morte\\
418: & Marc Magaña\\
419: & Marc Magaña\\
420: & Marc Magaña\\
421: & Juan J.~Rodríguez
\end{tabular}

\vspace{1em}

\noindent
Coordinador de les pràctiques: Alan Morte -- {alan.morte@uab.cat}

\noindent
Marc Magaña -- marc.magana@uab.cat

\noindent
Juan J.~Rodríguez -- {JuanJose.Rodriguez.Potosi@uab.cat
}


\subsubsection*{Avaluació}

\begin{itemize}
    \item La nota final de les pràctiques serà la mitjana de les dues proves d'avaluació.

    \item Les notes no són recuperables.

    \item Si hi ha més d'una falta injustificada, la nota final de les pràctiques serà un 0.
\end{itemize}

\subsubsection*{Procediment de les sessions}

\begin{itemize}
    \item Al campus virtual estan penjats els mòduls de les pràctiques. L'estudiant ha de preparar-se prèviament la pràctica amb el mòdul pertinent de la sessió.

    \item El professor està per ajudar i aclarir tots els dubtes que puguin sorgir. Però és important el treball autònom de l'estudiant i que segueixi les indicacions de la pràctica.

    \item A les sessions, el professor passarà llista d'assistència.

    \item Les faltes justificades han de ser avisades amb prou antelació i que hi hagi un motiu de pes. No es consideraran faltes justificades aquelles que no van ser avisades amb antelació (la setmana anterior, com a mínim). En aquest cas, si l'alumne ho demana, se li podrà reprogramar la sessió.

    En cas d'una urgència, s'haurà de presentar un document que l'avali.
\end{itemize}

\end{document}
